\chapter{Background and Related Work}
\label{ch:related}
% 개요 2장. 옛 비교(EVENT GRU, 소속 교환, history placement) 설명과
% 옛 2.4절 "Evidence granularity and forecasting baselines"(빈도 기준선)는
% 부록 B(app:secondary-results)로 옮긴다. 옮기는 작업은 부록 B를 쓸 때 한다.
% 옛 원고의 sec:recorded-case, sec:membership 참조는 이 장에서 뺐다.

\section{ATT\&CK, the kill chain, and Attack Flow}
\label{sec:related-attack}

MITRE ATT\&CK organizes observed adversary behavior into tactics, the
objectives an activity serves, and techniques and sub-techniques, the
ways an objective is pursued~\cite{mitreintro}; the Windows matrix lists
the techniques that apply to that platform~\cite{mitrewindows}. We use
this vocabulary to situate the behaviors in our data, but the 18
prediction labels (Appendix~\ref{app:actions}) are finer than ATT\&CK
techniques. They name operations at the level of the controlled
procedures. A system information enumeration and a registry query are
both discovery activity, and they relate to System Information Discovery
and Query Registry respectively; merging them into one output would erase
exactly the distinction the classifier is asked to make. Conversely,
different command lines that implement the same operation share a label.
The label is the operation designated by the procedure and verified
against the actual process start, not an ATT\&CK identifier.

The intrusion kill chain of Hutchins et al.\ groups a campaign into
phases so that defenders can reason about where to interrupt
it~\cite{hutchins2011killchain}. At a single host decision, however,
several operations fall under the same phase, and knowing that an actor
is in discovery does not say which resource it will examine next. Attack
Flow goes further and represents the order of and dependencies between
adversary actions, so a flow can document a known incident or a
hypothetical scenario~\cite{mitreattackflow}. Both are descriptions
written with the whole sequence in view. A forecaster has only the prefix
received before the next action, and it must commit to a distribution
over that action before it is observed.

\section{Attack graphs and attack progression}
\label{sec:related-graphs}

Classical attack graphs model how conditions and exploits combine to
reach a security-relevant state. Sheyner et al.\ generate such graphs
automatically from a system model~\cite{sheyner2002graphs}, and MulVAL
derives security consequences from configuration facts with logic
rules~\cite{ou2005mulval}. An edge in these graphs is a relation
defined by the security model, which lets an analyst compare alternative
routes and find the conditions several routes depend on. What the graph
covers is bounded by the supplied facts and rules.

These graphs answer the question of which paths are possible under a
known configuration. PROPHET asks a different one: given the part of an
execution that has become visible, how probable is each concrete next
action? A feasible path is informative for the first question even if
it never occurs, while the second requires probabilities that can be
scored against what the actor actually did; several feasible responses
to the same preceding action carry no information about which one this
actor will take. B\u{a}b\u{a}lau et al.\ move attack graphs toward the second
question by building them from intrusion-detection alerts and using
them to forecast attacker actions~\cite{babalau2024forecasting}.

\section{Forecasting attacker actions and security events}
\label{sec:related-forecasting}

Learned forecasters work from observed event streams rather than from a
configuration model. Tiresias trains a recurrent model on histories of
security events from intrusion-prevention telemetry and predicts the
next event~\cite{shen2018tiresias}. DeepCASE builds a context for each
security event from the events preceding it and uses attention over
that context to group events for analyst triage~\cite{vanede2022deepcase}.
In both, inputs and outputs are events at the granularity of the
monitored stream. Our outputs are instead semantic attacker actions at a
known boundary, and the records arriving before that action can include
the termination of a previous operation and unrelated benign work.

A second group forecasts at the level of techniques or graphs.
CL-AP$^{2}$ builds a time-aware attack knowledge graph, portrays past
behavior by tactic, and estimates a trajectory from which it derives the
adversary's next technique~\cite{liu2024clap2}. EFI extracts attack-scene
graphs from threat intelligence reports, matches them against endpoint
provenance, and forecasts the next interactions with a technique-level
interpretation~\cite{zhu2024efi}. BiTA combines bidirectional GRU
aggregation with a Transformer in a temporal graph network to predict
future network alerts~\cite{nayeri2026bita}. Their outputs are a
technique, an interaction graph, or an alert, whereas ours is a single
normalized distribution over 18 action categories.

Provenance-based detectors share the host-level evidence but not the
forecasting goal. KAIROS detects intrusions in whole-system provenance
graphs and reconstructs the activity relevant to an
investigation~\cite{cheng2024kairos}, and ORTHRUS detects and attributes
attacks at the level of provenance nodes~\cite{jiang2025orthrus}. Their
outputs concern behavior that has already been observed. Chapter~\ref{ch:generality}
runs adaptations of ORTHRUS, BiTA, and CL-AP$^{2}$ on our inputs as a
reference point only; EFI is not adapted. In the versions of the
systems in this section that we consulted, we did not find results
reported across random seeds or under benign workloads absent from
training. This statement is limited to those versions and is not a
claim that such evaluations do not exist elsewhere.
% 근거: RECORD:477, 480. 각 원문 판본에서 seed 변동 보고 여부를 사람이
% 다시 확인해 주세요(표 tab:related-tasks 마지막 열과 같은 범위).

\section{Sequence models for logs}
\label{sec:related-sequence}

Our inputs are Windows process creation and termination records, Security
events 4688 and 4689~\cite{ms4688, ms4689}, read only as they become
available at the receipt boundary defined in Chapter~\ref{ch:task}
(Figure~\ref{fig:cutoff}). Sequence models for system logs are the
closest general architectures for such input. DeepLog trains an LSTM
\cite{hochreiter1997long} on sequences of normal log keys and flags an entry as
anomalous when it falls outside the predicted next
keys~\cite{du2017deeplog}. LogBERT replaces the recurrent model with a
Transformer encoder~\cite{vaswani2017attention} trained by self-supervised
objectives on log sequences~\cite{guo2021logbert}. The gated recurrent unit
(GRU)~\cite{cho2014gru} is the simpler recurrent cell that our baseline
predictor uses at the event level (Chapter~\ref{ch:method}).

Chapter~\ref{ch:generality} uses three encoders drawn from this
family: an LSTM encoder (DeepLog-style), a Transformer encoder
(LogBERT-style), and a plain GRU. These are adaptations of the encoders
to our supervised next-action task, not reproductions of DeepLog or
LogBERT. They do not use the anomaly-detection objectives or
pretraining of the original systems; LogitNorm, input-gradient
regularization, and the relational residual of our baseline are
removed; and their settings were chosen by validation inside the FIT
set only.
% 근거: RECORD:281, 478.

\section{Underspecification and model agreement}
\label{sec:related-underspec}

D'Amour et al.\ describe underspecification: a training pipeline can
return many predictors with equivalent held-in performance that behave
differently under deployment conditions, and the pipeline does not
choose among them~\cite{damour2020underspecification}. Their work
establishes and documents the phenomenon across domains. We measure it
on next attacker action prediction and then try to reduce it. Baek et
al.\ observe that, across models, agreement on in-distribution data and
agreement on shifted data are linearly related in the same way as
accuracy, which allows performance under shift to be estimated from
unlabeled data~\cite{baek2022agreement}. That line of work uses agreement to
estimate performance; we use agreement as a training signal, letting a
student follow the consensus of several models on unlabeled environment
inputs. A link between agreement and performance should not be assumed
in our setting either: selecting seeds by agreement did not reproduce on
ADS-24, where the correlation was $+0.127$ ($p=0.503$;
Chapter~\ref{ch:ads24}).
% 근거: RECORD:405, 467--468.

\section{Distillation, ensembles, and unlabeled adaptation}
\label{sec:related-distillation}

Knowledge distillation trains a student to match the softened output
distribution of a teacher or of an ensemble~\cite{hinton2015distilling}.
Deep ensembles average the predictions of independently trained
networks and give better-calibrated predictive uncertainty than a single
network~\cite{lakshminarayanan2016simple}, at a cost that grows with the
number of members. Ensemble distribution distillation compresses an
ensemble into a single model that represents not only the mean
prediction but also the spread among members~\cite{malinin2019ensemble}. The
method in Chapter~\ref{ch:method} uses the simplest of these
ingredients: the arithmetic mean of selected teachers' probabilities
as the target, and a student with the teachers' architecture.

The comparison methods in Chapter~\ref{ch:generality} come from work
on stabilizing a single model or adapting it without labels.
Pseudo-labeling trains on the model's own most confident predictions
for unlabeled data~\cite{tmp_lee}, and entropy minimization penalizes
uncertain predictions on unlabeled inputs~\cite{tmp_grandvalet}.
Stochastic weight averaging (SWA) averages weights along one training
trajectory~\cite{izmailov2018averaging}, and temperature scaling rescales logits with
one parameter fitted on held-out data~\cite{guo2017calibration}. In security, Kan
et al.\ study adapting malware detectors to drift without labels,
where self-training on pseudo-labels can reinforce the model's own
errors~\cite{tmp_kan}. The distillation and adaptation work cited here
reports its gains as average performance on the target data. We seal
the seed standard deviation under unseen conditions as a primary
criterion, alongside the mean and the preservation of seen-condition
performance (Chapter~\ref{ch:evaluation}).
% 근거: RECORD:470--475.
% TODO(사람): 온도 보정과 확률 앙상블의 결과를 9.2절에 쓸지는 개요 질문 A1의
% 답에 따른다. 쓰지 않으면 이 문단의 temperature scaling 문장은 9.2절 기준선의
% 출처가 아니라 배경 설명으로만 남는다.

\section{Decomposition of ensemble log loss}
\label{sec:related-decomposition}

For the logarithmic loss, the natural way to combine probabilistic
predictors is the logarithmic opinion pool, the normalized geometric
mean of their distributions. Heskes showed that the log loss of such a
pool equals the average loss of its members minus an ambiguity term, the
average Kullback--Leibler divergence from the pool to the
members~\cite{heskes1997selecting}, and Gruber and Buettner place the same term
as the variance component in a bias--variance decomposition for proper
scores~\cite{gruber2022uncertainty}. Chapter~\ref{ch:decomposition} applies this
identity to the seeds of one configuration: mean seed loss splits
exactly into the loss of the log-pooled consensus and a disagreement
term $V$ that needs no labels. The identity is theirs; our use of it as
an unlabeled measure of instability in the deployment environment is a
post hoc analysis, introduced after the primary experiments.
% 근거: RECORD:308--316, 334, 469.

\section{Summary}
\label{sec:related-summary}

Table~\ref{tab:related-tasks} lists the evidence and output of each
work discussed above and whether variation across random seeds is
reported in the version we consulted.

\begin{table}[htbp]
\centering
\small
\caption{Evidence, outputs, and reported seed variation in related work.
The last column covers only the versions of the papers consulted for this
thesis.}
\label{tab:related-tasks}
\begin{tabularx}{\textwidth}{@{}p{0.20\textwidth}XXp{0.20\textwidth}@{}}
\toprule
Work & Evidence & Original output & Seed variation reported \\
\midrule
Sheyner et al.~\cite{sheyner2002graphs}
 & System model of hosts, vulnerabilities, and attacker
 & Attack graph of possible paths
 & Not applicable (no trained model) \\
MulVAL~\cite{ou2005mulval}
 & Configuration facts and logic rules
 & Derived security consequences
 & Not applicable (no trained model) \\
B\u{a}b\u{a}lau et al.~\cite{babalau2024forecasting}
 & Intrusion-detection alerts, as alert-driven attack graphs
 & Forecast of attacker actions
 & Not stated in the version consulted \\
Tiresias~\cite{shen2018tiresias}
 & Security event histories from intrusion-prevention telemetry
 & Next security event
 & Not stated in the version consulted \\
DeepCASE~\cite{vanede2022deepcase}
 & Sequences of security events
 & Event contexts grouped for analyst triage
 & Not stated in the version consulted \\
CL-AP$^{2}$~\cite{liu2024clap2}
 & Time-aware attack knowledge graph, tactic-level portrait
 & Next attack technique
 & Not stated in the version consulted \\
EFI~\cite{zhu2024efi}
 & Threat intelligence reports and endpoint provenance
 & Forecast interaction graph with technique-level interpretation
 & Not stated in the version consulted \\
BiTA~\cite{nayeri2026bita}
 & Temporal graph of network alerts
 & Future alerts
 & Not stated in the version consulted \\
KAIROS~\cite{cheng2024kairos}
 & Whole-system provenance graph
 & Intrusion detection and attack reconstruction
 & Not stated in the version consulted \\
ORTHRUS~\cite{jiang2025orthrus}
 & Whole-system provenance graph
 & Node-level detection and attribution
 & Not stated in the version consulted \\
DeepLog~\cite{du2017deeplog}
 & Sequences of system log keys
 & Anomaly flag from next-key prediction
 & Not stated in the version consulted \\
LogBERT~\cite{guo2021logbert}
 & Sequences of system log keys
 & Anomaly score from self-supervised training
 & Not stated in the version consulted \\
PROPHET (this thesis)
 & Windows audit records received before a known boundary
 & Distribution over 18 next attacker actions
 & 30 seeds per configuration; seed standard deviation is a sealed criterion \\
\bottomrule
\end{tabularx}
\end{table}
% TODO(사람): 마지막 열은 연구 기록(RECORD:477, 480)과 각 문헌의 초록 수준에서
% 정리했다. 원문 판본에서 seed 반복 보고 여부를 확인해 주세요. 보고가 있으면
% 해당 칸을 원문 표현으로 바꿔야 한다.
